\subsection{LLM Fallback}
\label{sec:llm-fallback}

Algorithm~\ref{alg:funspec-gen} presents the idealized symbolic-execution-driven path and intentionally omits the fallback layer; the LLM-fallback layer wraps that algorithm and is triggered in exactly two cases. (i) When the target function is self-recursive or manipulates recursive data structures: the symbolic executor neither computes recursive fixpoints nor enumerates the unbounded cell set of a recursive shape, so the fallback is invoked directly at the function entry instead of running the algorithm. (ii) When $\textit{SP}(\cdot)$ at the function exit returns $\textit{SymExecFailed}$: this covers the residual cases in which symbolic execution cannot produce a usable summary, including external calls, path explosion, and runtime errors. Loops and ordinary callees do not enter fallback; they are handled by $\mathrm{LoopInvGen}$ and recursive $\mathrm{FuncSpec}$ respectively. Their residual verifier failures are sent to the verifier-guided refinement layer, which also wraps Algorithm~\ref{alg:funspec-gen} rather than being inlined.

\rev{In the fallback, the LLM acts as a semantic summarization component for the missing symbolic-execution step: it receives the current code fragment, available symbolic states and path conditions, previously generated callee specifications or loop summaries, and an ACSL skeleton, and proposes ACSL clause candidates for four function-level tasks --- precondition generation, assigns-clause generation, postcondition generation, and inductive-predicate generation. The prompting inputs and the four tasks are detailed in \appref{app:fallback}.}

Fallback candidates are accepted only when Frama-C/WP proves them under the configured memory model and solver settings; verifier failures are handled by the refinement loop in Section~\ref{sec:refinement} (Table~\supref{tab:refine-strategy}, \appref{app:refine-ablation}). Fallback can therefore cover cases beyond the exact symbolic-execution fragment but 
unlike the postconditions computed by \(\textit{SP}\) (Proposition~\ref{thm:sp-sound-complete}), the results produced by the fallback are sound over‑approximations verified by the prover.

%does not extend the strongest-postcondition guarantee of Proposition~\ref{thm:sp-sound-complete}.
